\chapter{Problem Statement}

Most LMS platforms available today work the same way. Courses are organized as a flat list of topics, and student progress is measured through assignment scores, quiz marks, and attendance records. This gets the job done at a surface level, but it misses a lot.

\textbf{From the student's perspective,} grades are about the only feedback they get. A student can see that they scored 60\% on a quiz, but they cannot easily tell which topics connect to each other, which ones they missed, or what to study next. Progress information is scattered across tabs and pages. Getting a complete picture takes effort, and most students do not bother. The result is that many students feel lost in their courses without understanding exactly why.

\textbf{From the teacher's perspective,} a gradebook full of numbers does not say much about what is actually going on. It does not clearly show which students are falling behind, or which topics are confusing most of the class. By the time poor grades make the problem obvious, it is usually too late to intervene effectively.

\textbf{The gap between teaching and learning} is the core issue. A teacher moves through the syllabus at a planned pace. Not every student keeps up. Some fall behind on certain topics while doing fine on others. Current LMS tools have no good way to compare where the teacher is in the syllabus with where each student actually stands. This mismatch goes unnoticed until end-of-term exams make it impossible to ignore.

\textbf{Lack of visual engagement} is a separate but related problem. Cognitive science research shows that visual representations help people understand and retain information better \cite{eppler2006comparison}. Mind maps improve conceptual understanding. But most LMS platforms present everything as text and tables, which does not do much to hold a student's attention.

\textbf{AI integration concerns} make things worse. AI chatbots are becoming common in education, but most use general purpose language models that have no connection to any particular course. They can give answers that are technically correct but irrelevant to what the student is studying, or answers that are wrong entirely. Without being tied to the teacher's approved materials, these tools are unreliable.

The problem this research addresses is the absence of an LMS that can:
\begin{enumerate}
    \item Show students their learning progress as an interactive mindmap instead of a flat list of scores.
    \item Let teachers compare their teaching progress with each student's actual understanding in real time.
    \item Provide an AI study assistant that only uses the teacher's approved materials, reducing wrong or off-topic answers.
    \item Support multiple institutions on one platform without mixing their data.
\end{enumerate}
